%% Preambulo comun de los cuatro documentos del Subdocumento 13.

%% Legibilidad (BA Art. 40.4; Aclaraciones secciones 4 y 5)
%% 1. Sin expansion de fuente de microtype: con ella el cuerpo de 11 pt se
%%    compone a 10,8-10,9 pt y las tablas por debajo de su cuerpo nominal.
\microtypesetup{expansion=false}
%% 2. Titulos de tablas y figuras a 9,5 pt: \footnotesize da 8,9-9,0 pt,
%%    en el limite del minimo de 9 pt.
\DeclareCaptionFont{lfxcap}{\fontsize{9.5pt}{11.5pt}\selectfont}
\captionsetup{font={lfxcap,color=lafroxGris}}
%% 3. Sin palabras cortadas por guion dentro de las celdas (Aclaraciones,
%%    seccion 5: "sin palabras cortadas por columnas angostas"). La seccion 5
%%    rige solo para el cuerpo del subdocumento: el documento raiz del cuerpo
%%    activa \LFXsincortestrue; anexos y formularios conservan el corte.
\newif\ifLFXsincortes \LFXsincortesfalse
\renewcommand{\LFXcuerpotabla}{%
  \fontsize{9.5pt}{12.5pt}\selectfont
  \ifLFXsincortes\hyphenpenalty=10000\exhyphenpenalty=10000\fi}

%% Forma (BA Art. 40.2). Media firma del representante legal en el extremo
%% inferior derecho de cada pagina: filete de firma y rotulo sobre el folio.
%% La clase declara la opcion [firma] pero su pie derecho solo dibuja el
%% folio. Esta redefinicion debe llevarse a lafrox.cls en la rama base
%% (rama-formato-Latex-base) para que todos los subdocumentos la reciban.
%% La imagen de la firma se superpone al filete sin cambiar la altura del
%% pie: \raisebox[0pt][0pt] anula su alto y \rlap su ancho.
\newcommand{\LFXfirmaimagen}{logo/firma_representante.png}
\renewcommand{\LFXpiederecho}{%
  \begingroup\LFXpiebase
  \begin{tabular}[b]{@{}r@{}}
    {\rlap{\makebox[3.4cm][c]{\raisebox{-2pt}[0pt][0pt]{\includegraphics[height=18pt]{\LFXfirmaimagen}}}}%
     \color{lafroxGris}\rule{3.4cm}{0.4pt}}\\[-1pt]
    {\LFXchico Media firma: \LFXrepresentante}\\
    {\intermedium\fontsize{9pt}{11pt}\selectfont\color{lafroxTinta}\LFXfoliotexto}
  \end{tabular}%
  \endgroup}

%% Portada: raya tipografica real. El valor por defecto de la clase usa
%% "---", que la tipografia de la portada no convierte en raya.
\datoslafrox{sobre = {Sobre N.\textdegree{} 2 — Oferta Técnica}}

%% Tablas del cuerpo enteras en una pagina (Aclaraciones, seccion 5). En el
%% cuerpo (\LFXsincortestrue) todas las tablas caben en una pagina, asi que
%% tablalafrox se compone como flotante con tabularx: LaTeX la ubica donde
%% cabe, sin partirla y sin dejar huecos. flafter impide que una tabla o una
%% figura aparezca antes del punto del texto que la cita. Mismo estilo de la
%% clase: titulo arriba, \toprule redefinido con filetes entre filas.
\RequirePackage{flafter}
%% Una tabla debe aparecer antes del parrafo que interpreta sus resultados.
\RequirePackage{placeins}
%% Recorridos: impedir que el análisis posterior se adelante al diagrama.
\RequirePackage{float}
\newcommand{\recorridolafrox}[4]{%
  \begin{figure}[H]\centering
    \includegraphics[width=#1\textwidth]{#2}
    \caption{#3}\label{#4}
  \end{figure}}
\newcommand{\LFXtablasflotantes}{%
  \RenewDocumentEnvironment{tablalafrox}{m m m m +b}{%
    \begin{table}[htbp]
      \LFXcuerpotabla
      \caption{##1}\label{##2}
      \begin{tabularx}{\textwidth}{##3}
        \toprule ##4 \\ \midrule
        ##5
        \bottomrule
      \end{tabularx}
    \end{table}\FloatBarrier}{}}

%% tablasimple: tabla sin titulo (anexos y formularios), mismo estilo que
%% tablalafrox, con cabecera repetida y filetes entre filas.
\NewDocumentEnvironment{tablasimple}{m m +b}{%
  \par
  \LFXfiletear{#3}%
  \LFXcuerpotabla
  \setlength{\LTpre}{8pt}\setlength{\LTpost}{14pt}%
  \LFXconcuerpo{%
  \begin{xltabular}{\textwidth}{#1}
    \LFXtopruleoriginal #2 \\ \midrule
    \endfirsthead
    \LFXtopruleoriginal #2 \\ \midrule
    \endhead
    \midrule
    \endfoot
    \bottomrule
    \endlastfoot}%
}{}

%% URLs largas de las referencias: cortes en cualquier caracter.
\RequirePackage{xurl}

%% tablalafroxlarga: la tablalafrox original de la clase (xltabular con titulo,
%% cabecera repetida y "continuacion"). En el cuerpo, \LFXtablasflotantes
%% convierte tablalafrox en flotante de una pagina; esta variante se reserva a
%% las tablas que no caben en una pagina (declaracion de uso de IA).
\NewEnvironmentCopy{tablalafroxlarga}{tablalafrox}

%% Cifras y meses: alineacion derecha, con unidades en encabezado.
\newcolumntype{Z}{>{\raggedleft\arraybackslash}X}
